\spread{What a better choice of evidence changes}
\noindent The pictured question asks for the person who was beaten. Both names occur in the source. Changing how the reader scores a span changes the name delivered to the writer.
\platehere{Figure_171_reading_a_span_in_context}{Selecting the object of a verb need not preserve its role in the reply}{In panel a the revised reader selects the intended person. The writer uses that person as both the subject and object of the verb. In panel b a Chinese question asks when a bus service ended. Revising the selected span recovers the date, which suffices for a complete answer. Both examples show the original question, fact and stopped output images.}
\begin{multicols}{2}
\section*{Holding the writer fixed}
We compared the original reader with two learned scoring functions over its 32 highest-ranked positions. The control adds separate contributions from the two boundaries and the span length. The second model also sees the mean interior representation and combines both endpoints before applying a nonlinearity. These representations depend on the pictured question and page. Both functions add their correction to the original score. They have 70,144 trainable parameters each, start at the original score and see the same candidate lists in the same order.
\section*{Learning from ambiguous positions}
The objective sums the probabilities of all positions that match the annotated fact. For 29,299 of the 39,057 training questions at least one match was available. The remaining questions could supply no positive ranking gradient. If $m$ positions are correct, each has weight at most that of an incorrect winner whenever one is chosen. Their total weight is then at most $m$ times the combined rival weight. A Lean theorem checks this implication for nonnegative integer weights. On the held pages we checked it after exponentiating score differences and rounding at $10^{-9}$. For the whole-span scorer it certified two English and eighteen Chinese selections. This criterion concerns the chosen fact. The writer still has to express it.
\section*{Previously unused passages}
The new English cohort joins human QA2D statements to their source QAMR passages. We required a unique source sentence matching the question and containing the fact. The Chinese cohort uses CMRC. Each language contributes 96 separate passages and two distinct questions per passage. Original and revised readers evaluated the same page, question and candidate positions. The writer saw only the question and the chosen strip. Its output at the original first choice was reproduced on every question.
\end{multicols}
\newpage
\spread{The gap between knowing a name and writing a claim}
\noindent A changed date is often sufficient to answer a short question. A complete statement must also preserve who did what to whom.
\platehere{Figure_172_writing_after_span_selection}{Whole-span scoring improves Chinese short answers on the new passages}{The first two groups count a correct physical fact and a complete stopped answer image out of 192 questions per language. The third counts passages with both answers correct out of 96. The fourth counts both versions of an edited question, out of 52 English and 192 Chinese edits. The two added scoring functions have equal parameter counts and training exposure.}
\begin{multicols}{2}
\section*{Different outcomes in the two languages}
On Chinese questions, whole-span scoring recovered 17 facts lost by the original reader and lost five that it had selected correctly. Exact answer images increased from 31 to 43. The endpoint control gave 37. The 95\% passage-bootstrap intervals for the difference were 1.6--10.9 percentage points against the original and $-0.5$--7.3 against the endpoint control. Both answers were correct on six pages with whole-span scoring and three with the original. In English the three readers chose 11, 14 and 13 correct facts. All three writers completed the same five reference statements and solved zero natural pairs. The bilingual adoption criterion failed.
\section*{How much can reinserting the fact repair?}
We supplied the human fact on every page. The writer then produced 44 exact English and 187 exact Chinese answers. Among 143 English questions with a reference containing the fact exactly once, 113 complete images failed. Replacing any one contiguous interval of a generated reply with the correct fact could repair only 19. On 69 failures the fact was already present exactly once and on 15 it appeared more than once. On the 94 irreparable cases the reference also requires changes elsewhere in the reply. Some different wordings may be valid. A blinded local judge accepted only six of sixteen edited English human references and one of sixteen Chinese references, leaving semantic equivalence unresolved.
\section*{Generalization and cost}
On 32 questions the revised pipeline used 180 J of device energy. Q4 Qwen3-VL-8B used 9,240 J. Across sixteen English questions they included the requested fact once and twelve times, respectively. On 152 book prompts no pipeline reproduced the eight-glyph continuation. Across 130 Chinese historical-form requests the revised reader returned three exact specimens, with zero on 130 English requests. The study demonstrates limited transfer in reading short visual evidence. It leaves independent, general writing and matched-quality efficiency open.
\end{multicols}
\newpage
